\documentclass[10pt]{extarticle}
\usepackage[a4paper,margin=1.5cm]{geometry}
\usepackage[T1]{fontenc}
\usepackage{amsmath,amssymb,mathtools}
\usepackage{graphicx}
\usepackage[export]{adjustbox}
\usepackage{multicol}
\usepackage{enumitem}
\usepackage{booktabs,array,tabularx}
\setlength{\columnsep}{0.7cm}
\setlength{\parindent}{0pt}
\setlength{\parskip}{3pt}
\pagestyle{empty}
\makeatletter
\newcommand{\qfigdef}[4]{\@namedef{qf@file@#1}{#2}\@namedef{qf@w@#1}{#3}\@namedef{qf@wide@#1}{#4}}
\newcommand{\qf@check}[1]{\@ifundefined{qf@file@#1}{\PackageError{score2}{Unknown figure id #1}{}}{}}
\newcommand{\qf@img}[4][]{\includegraphics[width=#2\dimexpr\@nameuse{qf@w@#4}pt\relax,max width=#3,max height=0.45\textheight,#1]{\@nameuse{qf@file@#4}}}
\newif\ifqf@full
\newif\ifqf@boxed
\newsavebox{\qf@box}
\newcommand{\qf@open}{\global\qf@fulltrue\par\global\qf@boxedtrue
  \setbox\qf@box\vbox\bgroup\hsize\textwidth\linewidth\textwidth}
\newcommand{\qf@emit}[1]{\par\egroup\global\qf@boxedfalse
  \ifdim\dimexpr\ht\qf@box+\dp\qf@box\relax<#1\textheight\noindent\usebox{\qf@box}\par\else\unvbox\qf@box\fi}
\newcommand{\qfullsplit}{\ifqf@boxed\qf@emit{0.85}\fi}
\newcommand{\qf@close}{\ifqf@boxed\qf@emit{0.45}\fi\global\qf@fullfalse}
\newcommand{\qfullbegin}{\end{multicols}\qf@open}
\newcommand{\qfullend}{\qf@close\begin{multicols}{2}}
\newcommand{\qfullstart}{\qf@open}
\newcommand{\qfulldone}{\qf@close}
\newcommand{\qfig}[2][1]{\qf@check{#2}%
  \ifqf@full\qf@blk{#1}{#2}\else
  \ifnum\@nameuse{qf@wide@#2}=1 \expandafter\@firstoftwo\else\expandafter\@secondoftwo\fi
  {\end{multicols}\begin{center}\qf@img{#1}{\textwidth}{#2}\end{center}\begin{multicols}{2}}%
  {\qf@blk{#1}{#2}}\fi}
\newcommand{\qf@blk}[2]{\par\smallskip{\centering\qf@img{#1}{\linewidth}{#2}\par}\smallskip}
\newcommand{\qfigi}[2][1]{\qf@check{#2}\qf@img[valign=c]{#1}{\linewidth}{#2}}
\makeatother
\newcommand{\papersec}[1]{\par\bigskip{\centering\Large\bfseries PAPER-#1\par}\medskip}
\newcommand{\partsec}[2]{\par\medskip\noindent{\large\bfseries PART-#1: #2}\par\smallskip}
\newcommand{\sect}[1]{\par\smallskip\noindent\textbf{#1}\par}
\newcommand{\qn}[1]{\par\smallskip\noindent\textbf{#1.}~}
\newcommand{\paraq}[2]{\par\smallskip\noindent\textbf{Paragraph for Questions #1 and #2}\par}
\newenvironment{qoptions}{\begin{itemize}[leftmargin=1.8em,itemsep=1pt,topsep=2pt,label={}]}{\end{itemize}}
\newcommand{\opt}[2]{\item[(#1)] #2}
\newcolumntype{L}[1]{>{\raggedright\arraybackslash}p{#1}}
\newcommand{\fwstart}{\end{multicols}\par\noindent\begin{minipage}{\textwidth}\small\setlength{\tabcolsep}{4pt}\renewcommand{\arraystretch}{1.5}}
\newcommand{\fwend}{\end{minipage}\par\medskip\begin{multicols}{2}}
% figure registry (generated by score2_tools.py finalize)
\qfigdef{p1_phy_s1_q01_1}{SRG (Enthusiast) - Test 11 (29 Aug 2021)-img/p1_phy_s1_q01_1.png}{106.7}{0}
\qfigdef{p1_phy_s1_q02_1}{SRG (Enthusiast) - Test 11 (29 Aug 2021)-img/p1_phy_s1_q02_1.png}{98.4}{0}
\qfigdef{p1_phy_s1_q03_1}{SRG (Enthusiast) - Test 11 (29 Aug 2021)-img/p1_phy_s1_q03_1.png}{121.7}{0}
\qfigdef{p1_phy_s1_q04_1}{SRG (Enthusiast) - Test 11 (29 Aug 2021)-img/p1_phy_s1_q04_1.png}{119.8}{0}
\qfigdef{p1_phy_s1_q06_1}{SRG (Enthusiast) - Test 11 (29 Aug 2021)-img/p1_phy_s1_q06_1.png}{144.6}{0}
\qfigdef{p1_phy_s2_q09_1}{SRG (Enthusiast) - Test 11 (29 Aug 2021)-img/p1_phy_s2_q09_1.png}{76.8}{0}
\qfigdef{p1_phy_s2_q10_1}{SRG (Enthusiast) - Test 11 (29 Aug 2021)-img/p1_phy_s2_q10_1.png}{68.5}{0}
\qfigdef{p1_phy_s2_q11_1}{SRG (Enthusiast) - Test 11 (29 Aug 2021)-img/p1_phy_s2_q11_1.png}{95.1}{0}
\qfigdef{p1_phy_s3_q14_1}{SRG (Enthusiast) - Test 11 (29 Aug 2021)-img/p1_phy_s3_q14_1.png}{197.5}{0}
\qfigdef{p1_phy_s3_q15_1}{SRG (Enthusiast) - Test 11 (29 Aug 2021)-img/p1_phy_s3_q15_1.png}{51.5}{0}
\qfigdef{p1_phy_s3_q17_1}{SRG (Enthusiast) - Test 11 (29 Aug 2021)-img/p1_phy_s3_q17_1.png}{123.3}{0}
\qfigdef{p1_phy_s3_q18_1}{SRG (Enthusiast) - Test 11 (29 Aug 2021)-img/p1_phy_s3_q18_1.png}{112.8}{0}
\qfigdef{p1_chem_s1_q01_1}{SRG (Enthusiast) - Test 11 (29 Aug 2021)-img/p1_chem_s1_q01_1.png}{245.2}{0}
\qfigdef{p1_chem_s1_q01_A}{SRG (Enthusiast) - Test 11 (29 Aug 2021)-img/p1_chem_s1_q01_A.png}{82.9}{0}
\qfigdef{p1_chem_s1_q01_B}{SRG (Enthusiast) - Test 11 (29 Aug 2021)-img/p1_chem_s1_q01_B.png}{83.4}{0}
\qfigdef{p1_chem_s1_q01_C}{SRG (Enthusiast) - Test 11 (29 Aug 2021)-img/p1_chem_s1_q01_C.png}{82.9}{0}
\qfigdef{p1_chem_s1_q01_D}{SRG (Enthusiast) - Test 11 (29 Aug 2021)-img/p1_chem_s1_q01_D.png}{68.3}{0}
\qfigdef{p1_chem_s1_q01_6}{SRG (Enthusiast) - Test 11 (29 Aug 2021)-img/p1_chem_s1_q01_6.png}{174.5}{0}
\qfigdef{p1_chem_s1_q02_A}{SRG (Enthusiast) - Test 11 (29 Aug 2021)-img/p1_chem_s1_q02_A.png}{154.7}{0}
\qfigdef{p1_chem_s1_q02_B}{SRG (Enthusiast) - Test 11 (29 Aug 2021)-img/p1_chem_s1_q02_B.png}{125.3}{0}
\qfigdef{p1_chem_s1_q02_D}{SRG (Enthusiast) - Test 11 (29 Aug 2021)-img/p1_chem_s1_q02_D.png}{111.1}{0}
\qfigdef{p1_chem_s1_q03_1}{SRG (Enthusiast) - Test 11 (29 Aug 2021)-img/p1_chem_s1_q03_1.png}{227.1}{0}
\qfigdef{p1_chem_s1_q03_A}{SRG (Enthusiast) - Test 11 (29 Aug 2021)-img/p1_chem_s1_q03_A.png}{104.1}{0}
\qfigdef{p1_chem_s1_q03_B}{SRG (Enthusiast) - Test 11 (29 Aug 2021)-img/p1_chem_s1_q03_B.png}{95.4}{0}
\qfigdef{p1_chem_s1_q03_C}{SRG (Enthusiast) - Test 11 (29 Aug 2021)-img/p1_chem_s1_q03_C.png}{188.1}{0}
\qfigdef{p1_chem_s1_q03_D}{SRG (Enthusiast) - Test 11 (29 Aug 2021)-img/p1_chem_s1_q03_D.png}{91.8}{0}
\qfigdef{p1_chem_s1_q03_5}{SRG (Enthusiast) - Test 11 (29 Aug 2021)-img/p1_chem_s1_q03_5.png}{53.0}{0}
\qfigdef{p1_chem_s1_q04_1}{SRG (Enthusiast) - Test 11 (29 Aug 2021)-img/p1_chem_s1_q04_1.png}{178.9}{0}
\qfigdef{p1_chem_s1_q04_A}{SRG (Enthusiast) - Test 11 (29 Aug 2021)-img/p1_chem_s1_q04_A.png}{50.0}{0}
\qfigdef{p1_chem_s1_q04_B}{SRG (Enthusiast) - Test 11 (29 Aug 2021)-img/p1_chem_s1_q04_B.png}{63.5}{0}
\qfigdef{p1_chem_s1_q04_C}{SRG (Enthusiast) - Test 11 (29 Aug 2021)-img/p1_chem_s1_q04_C.png}{57.6}{0}
\qfigdef{p1_chem_s1_q04_D}{SRG (Enthusiast) - Test 11 (29 Aug 2021)-img/p1_chem_s1_q04_D.png}{50.2}{0}
\qfigdef{p1_chem_s1_q05_1}{SRG (Enthusiast) - Test 11 (29 Aug 2021)-img/p1_chem_s1_q05_1.png}{164.3}{0}
\qfigdef{p1_chem_s1_q05_2}{SRG (Enthusiast) - Test 11 (29 Aug 2021)-img/p1_chem_s1_q05_2.png}{68.7}{0}
\qfigdef{p1_chem_s1_q05_3}{SRG (Enthusiast) - Test 11 (29 Aug 2021)-img/p1_chem_s1_q05_3.png}{72.2}{0}
\qfigdef{p1_chem_s1_q05_5}{SRG (Enthusiast) - Test 11 (29 Aug 2021)-img/p1_chem_s1_q05_5.png}{82.3}{0}
\qfigdef{p1_chem_s1_q05_4}{SRG (Enthusiast) - Test 11 (29 Aug 2021)-img/p1_chem_s1_q05_4.png}{72.2}{0}
\qfigdef{p1_chem_s1_q05_6}{SRG (Enthusiast) - Test 11 (29 Aug 2021)-img/p1_chem_s1_q05_6.png}{72.5}{0}
\qfigdef{p1_chem_s1_q05_7}{SRG (Enthusiast) - Test 11 (29 Aug 2021)-img/p1_chem_s1_q05_7.png}{82.5}{0}
\qfigdef{p1_chem_s1_q05_8}{SRG (Enthusiast) - Test 11 (29 Aug 2021)-img/p1_chem_s1_q05_8.png}{82.5}{0}
\qfigdef{p1_chem_s1_q05_9}{SRG (Enthusiast) - Test 11 (29 Aug 2021)-img/p1_chem_s1_q05_9.png}{72.2}{0}
\qfigdef{p1_chem_s2_q07_1}{SRG (Enthusiast) - Test 11 (29 Aug 2021)-img/p1_chem_s2_q07_1.png}{39.5}{0}
\qfigdef{p1_chem_s2_q07_A}{SRG (Enthusiast) - Test 11 (29 Aug 2021)-img/p1_chem_s2_q07_A.png}{70.9}{0}
\qfigdef{p1_chem_s2_q07_B}{SRG (Enthusiast) - Test 11 (29 Aug 2021)-img/p1_chem_s2_q07_B.png}{52.8}{0}
\qfigdef{p1_chem_s2_q07_C}{SRG (Enthusiast) - Test 11 (29 Aug 2021)-img/p1_chem_s2_q07_C.png}{66.5}{0}
\qfigdef{p1_chem_s2_q07_D}{SRG (Enthusiast) - Test 11 (29 Aug 2021)-img/p1_chem_s2_q07_D.png}{48.0}{0}
\qfigdef{p1_chem_s2_q08_1}{SRG (Enthusiast) - Test 11 (29 Aug 2021)-img/p1_chem_s2_q08_1.png}{267.3}{1}
\qfigdef{p1_chem_s2_q08_3}{SRG (Enthusiast) - Test 11 (29 Aug 2021)-img/p1_chem_s2_q08_3.png}{85.5}{0}
\qfigdef{p1_chem_s2_q08_B}{SRG (Enthusiast) - Test 11 (29 Aug 2021)-img/p1_chem_s2_q08_B.png}{126.5}{0}
\qfigdef{p1_chem_s2_q08_C}{SRG (Enthusiast) - Test 11 (29 Aug 2021)-img/p1_chem_s2_q08_C.png}{85.5}{0}
\qfigdef{p1_chem_s2_q08_5}{SRG (Enthusiast) - Test 11 (29 Aug 2021)-img/p1_chem_s2_q08_5.png}{107.1}{0}
\qfigdef{p1_chem_s2_q09_A}{SRG (Enthusiast) - Test 11 (29 Aug 2021)-img/p1_chem_s2_q09_A.png}{144.2}{0}
\qfigdef{p1_chem_s2_q09_B}{SRG (Enthusiast) - Test 11 (29 Aug 2021)-img/p1_chem_s2_q09_B.png}{134.6}{0}
\qfigdef{p1_chem_s2_q09_D}{SRG (Enthusiast) - Test 11 (29 Aug 2021)-img/p1_chem_s2_q09_D.png}{155.1}{0}
\qfigdef{p1_chem_s2_q11_A}{SRG (Enthusiast) - Test 11 (29 Aug 2021)-img/p1_chem_s2_q11_A.png}{140.7}{0}
\qfigdef{p1_chem_s2_q11_B}{SRG (Enthusiast) - Test 11 (29 Aug 2021)-img/p1_chem_s2_q11_B.png}{127.2}{0}
\qfigdef{p1_chem_s2_q11_C}{SRG (Enthusiast) - Test 11 (29 Aug 2021)-img/p1_chem_s2_q11_C.png}{116.3}{0}
\qfigdef{p1_chem_s2_q11_D}{SRG (Enthusiast) - Test 11 (29 Aug 2021)-img/p1_chem_s2_q11_D.png}{126.7}{0}
\qfigdef{p1_chem_s2_q12_S}{SRG (Enthusiast) - Test 11 (29 Aug 2021)-img/p1_chem_s2_q12_S.png}{119.1}{0}
\qfigdef{p1_chem_s2_q12_A}{SRG (Enthusiast) - Test 11 (29 Aug 2021)-img/p1_chem_s2_q12_A.png}{101.0}{0}
\qfigdef{p1_chem_s2_q12_B}{SRG (Enthusiast) - Test 11 (29 Aug 2021)-img/p1_chem_s2_q12_B.png}{118.5}{0}
\qfigdef{p1_chem_s2_q12_C}{SRG (Enthusiast) - Test 11 (29 Aug 2021)-img/p1_chem_s2_q12_C.png}{150.7}{0}
\qfigdef{p1_chem_s2_q12_D}{SRG (Enthusiast) - Test 11 (29 Aug 2021)-img/p1_chem_s2_q12_D.png}{150.7}{0}
\qfigdef{p1_chem_s3_q13_1}{SRG (Enthusiast) - Test 11 (29 Aug 2021)-img/p1_chem_s3_q13_1.png}{375.1}{1}
\qfigdef{p1_chem_s3_q15_1}{SRG (Enthusiast) - Test 11 (29 Aug 2021)-img/p1_chem_s3_q15_1.png}{175.5}{0}
\qfigdef{p1_chem_s3_q17_1}{SRG (Enthusiast) - Test 11 (29 Aug 2021)-img/p1_chem_s3_q17_1.png}{166.5}{0}
\qfigdef{p1_chem_s3_q18_1}{SRG (Enthusiast) - Test 11 (29 Aug 2021)-img/p1_chem_s3_q18_1.png}{183.7}{0}
\raggedcolumns\makeatletter\@beginparpenalty=10000\makeatother

\begin{document}
\begin{center}{\Large\bfseries SRG (Enthusiast) --- Test 11 (29 Aug 2021)}\end{center}
\begin{multicols}{2}
\partsec{1}{PHYSICS}
\sect{SECTION-I (i)}
\qn{1} A small coil $C$ with $N=200$ turns is mounted on one end of a balance beam and introduced between the poles of an electromagnet as shown in the figure. The area of the coil is $S=1\ \text{cm}^2$, then length of the right arm of the balance beam is $l=30$ cm. When there is no current in the coil the balance is in equilibrium. On passing a current $I=22$ mA through the coil, equilibrium is restored by putting an additional weight of mass $m=60$ mg on the balance pan. Magnetic field between poles (Assume constant)
\qfig{p1_phy_s1_q01_1}
\begin{qoptions}
\opt{A}{0.8 Tesla}
\opt{B}{0.4 Tesla}
\opt{C}{0.2 Tesla}
\opt{D}{1 Tesla}
\end{qoptions}
\qn{2} In figure there is a uniform conducting structure in which each small square has side $a$. The structure is kept in a uniform magnetic field $B$. Then the magnetic force on the structure will be.
\qfig{p1_phy_s1_q02_1}
\begin{qoptions}
\opt{A}{$2\sqrt2\,iBa$}
\opt{B}{$\sqrt2\,iBa$}
\opt{C}{$2\,iBa$}
\opt{D}{$iBa$}
\end{qoptions}
\qn{3} A light rod of length 4 m is connected rigidly with two identical particles each of mass $m=2$ kg. The free end of the rod is smoothly pivoted. The rod is released from rest from its horizontal position. Reaction force (in N) offered by the pivot just after release. ($g=10\ \text{m/s}^2$) is
\qfig{p1_phy_s1_q03_1}
\begin{qoptions}
\opt{A}{2N}
\opt{B}{4N}
\opt{C}{6N}
\opt{D}{8N}
\end{qoptions}
\qn{4} A wheel $A$ is connected to a second wheel $B$ by means of an inextensible light string passing over a pulley $C$, which rotates about a fixed horizontal axle $O$ as shown in the figure. The system is released from rest. The wheel $A$ rolls down the inclined plane $OK$ thus pulling up wheel $B$ which rolls up along the inclined plane $ON$. Determine the velocity (in m/s) of the axle of wheel $A$, when it has travelled a distance $s=3.5$ m down the slope. Both wheels and pulley are assumed to be homogeneous disks of identical weight and radius. The string does not slip over $C$. (take $\alpha=53^\circ$ and $\beta=37^\circ$, $g=10\ \text{m/s}^2$)
\qfig{p1_phy_s1_q04_1}
\begin{qoptions}
\opt{A}{1 m/s}
\opt{B}{2 m/s}
\opt{C}{4 m/s}
\opt{D}{8 m/s}
\end{qoptions}
\qn{5} A solid cylinder is released from rest from the top of an incline plane of inclination $60^\circ$ where friction coefficient varies with distance $x$ as $\mu=\dfrac{2-3x}{\sqrt3}$. Find the distance travelled by the cylinder on incline before it starts slipping:
\begin{qoptions}
\opt{A}{$1/3$ m}
\opt{B}{$1/\sqrt3$ m}
\opt{C}{3 m}
\opt{D}{$\sqrt3$ m}
\end{qoptions}
\qn{6} A metal wire $PQ$ of mass 10g lies at rest on two horizontal metal rails separated by 5 cm as shown in the figure. A vertically downward magnetic field of magnitude 0.80 T exists in the space. The resistance of the circuit is slowly decreased and it is found that when the resistance goes below $20.0\Omega$, the wire $PQ$ starts sliding on the rails. The coefficient of friction between wire and rail is (Acceleration due to gravity $=10\ \text{ms}^{-2}$)
\qfig{p1_phy_s1_q06_1}
\begin{qoptions}
\opt{A}{3/25}
\opt{B}{9/25}
\opt{C}{6/25}
\opt{D}{2}
\end{qoptions}
\sect{SECTION-I (ii)}
\qn{7} A cyclist goes around a curve of radius $r$ on a horizontal road with speed $v$. Which of the following options are correct?
\begin{qoptions}
\opt{A}{The required centripetal acceleration is provided by component of normal reaction}
\opt{B}{The required centripetal acceleration is provided by friction}
\opt{C}{The net force by the road on the cyclist should pass through the centre of mass of the system of cycle and cyclist}
\opt{D}{The angle made by the cyclist with the vertical while going around the curve is $\tan^{-1}\left(\dfrac{v^2}{rg}\right)$}
\end{qoptions}
\qn{8} A positive charge $(q,m)$ moving velocity $v_0\hat i$ along $x$-axis enters a region of magnetic field $\vec B=B_0x\hat k$ ($B_0$ is a positive constant and $x\ge0$). Choose the correct statements regarding the subsequent motion of the particle.
\begin{qoptions}
\opt{A}{It moves on a circular path}
\opt{B}{It moves on a non-circular path}
\opt{C}{Radius of curvature of path when it crosses plane $x=L(L>0)$ is $\dfrac{mv_0}{qB_0L}$}
\opt{D}{Radius of curvature of path when it crosses plane $x=L(L>0)$ is $\dfrac{mv_0}{qB_0}$}
\end{qoptions}
\qn{9} A thin uniform rod of mass $m$ and length $l$ is hinged at one end. It is released from horizontal position as shown in figure. The hinge is smooth. Choose the correct options.
\qfig{p1_phy_s2_q09_1}
\begin{qoptions}
\opt{A}{The initial angular acceleration of the rod is $\dfrac{3g}{2\ell}$}
\opt{B}{The force exerted by hinge on the rod just after release is $\dfrac{mg}{4}$}
\opt{C}{The angular velocity of the rod when it becomes vertical is $\sqrt{\dfrac{3g}{\ell}}$}
\opt{D}{The force exerted by hinge on the rod when it becomes vertical is $mg$}
\end{qoptions}
\qn{10} A circular conducting loop of radius $r_0$ and having resistance per unit length $\lambda$ as shown in the figure is placed in a magnetic field $B$ which is constant in space and time. The ends of the loop are crossed and pulled in opposite directions with a velocity $v$ such that the loop always remains circular and the radius of the loop goes on decreasing, then choose the correct options.
\qfig{p1_phy_s2_q10_1}
\begin{qoptions}
\opt{A}{Radius of the loop changes with times as $r=r_0-\dfrac{vt}{\pi}$}
\opt{B}{EMF induced in the loop as a function time is $\varepsilon=2Bv\left[r_0-\dfrac{vt}{\pi}\right]$}
\opt{C}{Current induced in the loop is $I=\dfrac{Bv}{2\pi\lambda}$}
\opt{D}{Current induced in the loop is $I=\dfrac{Bv}{\pi\lambda}$}
\end{qoptions}
\qn{11} A solid sphere $A$ is rolling without slipping on a rough surface with angular speed $\omega$ as shown in figure. It collides elastically with another identical sphere $B$ at rest. Assume no friction acts between the spheres during collision. Radius of the spheres is $R$, then choose correct option/s
\qfig{p1_phy_s2_q11_1}
\begin{qoptions}
\opt{A}{$A$ and $B$ move in opposite directions}
\opt{B}{$A$ and $B$ move in same direction}
\opt{C}{A long time after the collision, the ratio of velocity of centres of the spheres $A$ and $B$ is 0.4}
\opt{D}{A long time after the collision, the ratio of velocity of centres of the spheres $A$ and $B$ is 0.67}
\end{qoptions}
\sect{SECTION-II}
\qn{13} The moment of inertia of a uniform hollow cubical box of mass $M$ and side $a$ about an axis passing through the centre of two opposite faces is equal to $xMa^2$, where $x$ is \underline{\hspace{1.5em}}
\qn{14} A segment of length $2a$ is cut along $Z$ axis from an infinite sheet having a surface current per unit width $J$ coming out of the plane in the given diagram. The magnetic field at $P$ exactly above the midpoint of the cut at a height $h=\dfrac{a}{\sqrt3}$ is $\dfrac{\mu_0J}{n}(-\hat i)$. Find the value of $n$.
\qfig{p1_phy_s3_q14_1}
\qn{15} A thin uniform rod of mass $2m$ and length $L$, hinged at its upper end, hangs in equilibrium. A small bullet of mass $m$ is fired horizontally with velocity $v$. The bullet hits the rod very close to its bottom end and gets embedded in it. The angular velocity of the combined system just afterwards is $X\left(\dfrac vL\right)$. The value of $X$ is \underline{\hspace{2em}}.
\qfig{p1_phy_s3_q15_1}
\qn{16} A charged particle is projected in magnetic field $\vec B=10\hat k$ from origin in $x-y$ plane. The Particle moves in a circle and just touches the line $y=5$m at $x=5\sqrt3$m. Find the speed of particle (in m/s) given that its mass $=5\times10^{-5}$ kg and charge $=1\mu$C
\qn{17} A solid cylinder of mass 10 kg and diameter 20 cm mounted to rotate about its central axis is initially at rest. A torque begins to act on the cylinder at $t=0$. The variation with time of this torque is given in the graph. The angular speed of the cylinder at $t=20$ sec is \underline{\hspace{1.5em}} rad/s.
\qfig{p1_phy_s3_q17_1}
\qn{18} Given figure shows an amperian loop $ABCDA$. Part $ABC$ is in vertical plane $PSTU$ while part $CDA$ is in horizontal plane $PQRS$. Direction of circulation along the path is shown by an arrow near point $B$ and $D$. $\oint\vec B\cdot d\vec l$ for this path according Ampere's law will be $\mu_0(ai_1+bi_2+ci_3)$. Then mark your answer $a+b+c$.
\qfig{p1_phy_s3_q18_1}
\partsec{2}{CHEMISTRY}
\sect{SECTION-I (i)}
\qn{1}
\qfig{p1_chem_s1_q01_1}
The product is
\begin{qoptions}
\opt{A}{\qfigi{p1_chem_s1_q01_A}}
\opt{B}{\qfigi{p1_chem_s1_q01_B}}
\opt{C}{\qfigi{p1_chem_s1_q01_C}}
\opt{D}{\qfigi{p1_chem_s1_q01_D}}
\end{qoptions}
\columnbreak\qn{2}
\qfig{p1_chem_s1_q01_6}
$B$ is
\begin{qoptions}
\opt{A}{\qfigi{p1_chem_s1_q02_A}}
\opt{B}{\qfigi{p1_chem_s1_q02_B}}
\opt{C}{Both (A) and (B)}
\opt{D}{\qfigi{p1_chem_s1_q02_D}}
\end{qoptions}
\qn{3}
\qfig{p1_chem_s1_q03_1}
Where $X$ is a compound that forms azo dye with benzene diazonium chloride in faintly basic medium. Hence the products $P$, $X$ and $Y$ are respectively
\begin{qoptions}
\opt{A}{\qfigi{p1_chem_s1_q03_A} and $\mathrm{PhCOCH_3}$}
\opt{B}{\qfigi{p1_chem_s1_q03_B} and $\mathrm{PhCH_2CHO}$}
\opt{C}{\qfigi{p1_chem_s1_q03_C}}
\opt{D}{\qfigi{p1_chem_s1_q03_D} and \qfigi{p1_chem_s1_q03_5}}
\end{qoptions}
\qn{4}
\qfig{p1_chem_s1_q04_1}
\begin{qoptions}
\opt{A}{\qfigi{p1_chem_s1_q04_A}}
\opt{B}{\qfigi{p1_chem_s1_q04_B}}
\opt{C}{\qfigi{p1_chem_s1_q04_C}}
\opt{D}{\qfigi{p1_chem_s1_q04_D}}
\end{qoptions}
\qn{5}
\qfig{p1_chem_s1_q05_1}
\begin{qoptions}
\opt{A}{$P$: \qfigi{p1_chem_s1_q05_2}\quad $Q$: \qfigi{p1_chem_s1_q05_3}}
\opt{B}{$P$: \qfigi{p1_chem_s1_q05_5}\quad $Q$: \qfigi{p1_chem_s1_q05_4}}
\opt{C}{$P$: \qfigi{p1_chem_s1_q05_6}\quad $Q$: \qfigi{p1_chem_s1_q05_7}}
\opt{D}{$P$: \qfigi{p1_chem_s1_q05_8}\quad $Q$: \qfigi{p1_chem_s1_q05_9}}
\end{qoptions}
\qn{6} In a system $\mathrm{A_{(s)}\rightleftharpoons2B_{(g)}+3C_{(g)}}$. If the concentration of C at equilibrium is increased by a factor 2, it will cause the equilibrium concentration of B to change to
\begin{qoptions}
\opt{A}{Two times of its original value}
\opt{B}{One half its original value}
\opt{C}{$2\sqrt2$ time of its original value}
\opt{D}{$\dfrac1{2\sqrt2}$ times of its original value}
\end{qoptions}
\sect{SECTION-I (ii)}
\qn{7} Which of the following are formed during Reimer - Tiemann reaction of the compound
\qfig{p1_chem_s2_q07_1}
\begin{qoptions}
\opt{A}{\qfigi{p1_chem_s2_q07_A}}
\opt{B}{\qfigi{p1_chem_s2_q07_B}}
\opt{C}{\qfigi{p1_chem_s2_q07_C}}
\opt{D}{\qfigi{p1_chem_s2_q07_D}}
\end{qoptions}
\qfullbegin
\qn{8} Which of the following are correct in the given sequence of reaction
\qfig{p1_chem_s2_q08_1}
\qfullsplit
\begin{qoptions}
\opt{A}{\qfigi{p1_chem_s2_q08_3}}
\opt{B}{\qfigi{p1_chem_s2_q08_B}}
\opt{C}{\qfigi{p1_chem_s2_q08_C}}
\opt{D}{\qfigi{p1_chem_s2_q08_5}}
\end{qoptions}
\qfullend
\qn{9} Identify reactions give ketone product?
\begin{qoptions}
\opt{A}{\qfigi{p1_chem_s2_q09_A}}
\opt{B}{\qfigi{p1_chem_s2_q09_B}}
\opt{C}{$\mathrm{PhMgBr}\xrightarrow[(2)\,\mathrm{H_3O^+}]{(1)\,\mathrm{Me-C\equiv N}}$}
\opt{D}{\qfigi{p1_chem_s2_q09_D}}
\end{qoptions}
\qn{10} $\underset{\text{(solid)}}{\mathrm{NH_2COONH_4}}\rightleftharpoons2\underset{\text{(gas)}}{\mathrm{NH_3}}+\underset{\text{(gas)}}{\mathrm{CO_2}}$

At equilibrium state, which of the following features are correct?
\begin{qoptions}
\opt{A}{Total pressure is 1.5 timed that of partial pressure of $\mathrm{NH_3}$.}
\opt{B}{$K_p$ value is 4 times that of $\left(P_{\mathrm{CO_2}}\right)^3$.}
\opt{C}{$P_{\mathrm{NH_3}}=\dfrac12P_{\mathrm{CO_2}}$}
\opt{D}{$\mathrm{CO_2}$ contributes 33.3\% of equilibrium pressure}
\end{qoptions}
\qn{11} In which reactions the product will be ketone?
\begin{qoptions}
\opt{A}{\qfigi{p1_chem_s2_q11_A}}
\opt{B}{\qfigi{p1_chem_s2_q11_B}}
\opt{C}{\qfigi{p1_chem_s2_q11_C}}
\opt{D}{\qfigi{p1_chem_s2_q11_D}}
\end{qoptions}
\qn{12} If ester $A$ contain five carbon atom and product $B$ and $C$ gives positive iodoform test then structure of ester $A$ is / are
\qfig{p1_chem_s2_q12_S}
\begin{qoptions}
\opt{A}{\qfigi{p1_chem_s2_q12_A}}
\opt{B}{\qfigi{p1_chem_s2_q12_B}}
\opt{C}{\qfigi{p1_chem_s2_q12_C}}
\opt{D}{\qfigi{p1_chem_s2_q12_D}}
\end{qoptions}
\qfullbegin
\sect{SECTION-II}
\qn{13} How many ether will give two alkyl iodides on reaction with excess amount of HI
\qfig{p1_chem_s3_q13_1}
\qfullend
\qn{14} A Compound $A$ with formula $\mathrm{C_4H_{10}O_3}$ on acylation with acetic anhydride gives another compound $B$ with molecular mass 190. The number of hydroxyl group in the compound $A$ are
\qn{15} When 3.06 g of solid $\mathrm{NH_4HS}$ is introduced into a two-liter evacuated flask at 27$^\circ$C, 30\% of the solid decomposes into gaseous ammonia and hydrogen sulphide

Calculate $K_P$ for the reaction at 27$^\circ$C
\qn{16}
\qfig{p1_chem_s3_q15_1}
What is value of $x$
\qn{17}
\qfig{p1_chem_s3_q17_1}
Number of $\mathrm O$ atoms in $B$ are $x$

Number of $\mathrm O$ atoms in $C$ are $y$

then $x+y=$ \underline{\hspace{1.5em}}
\qn{18}
\qfig{p1_chem_s3_q18_1}
Degree of unsaturation of $D$.
\partsec{3}{MATHEMATICS}
\sect{SECTION-I (i)}
\qn{1} The integral $\displaystyle\int\dfrac{x^3}{(1-x\cot x)^2}dx$ is
\begin{qoptions}
\opt{A}{$\dfrac{-x^3}{\tan x-x}-\dfrac{3x^2}2+3\log|\sin x-x\cos x|+c$}
\opt{B}{$\dfrac{-x^3}{\tan x-x}+\dfrac{x^2}2+3\log|\sin x-x\cos x|+c$}
\opt{C}{$\dfrac{x^3}{\tan x-x}-\dfrac{3x^2}2+3\log|\sin x-x\cos x|+c$}
\opt{D}{$\dfrac{-x^3}{\tan x-x}-\dfrac{x^2}2+3\log|\sin x-x\cos x|+c$}
\end{qoptions}
\qn{2} If $\displaystyle\int x^{26}(x-1)^{17}(5x-3)dx=\dfrac{x^{27}(x-1)^{18}}{k}+C$, where $C$ is constant of integration, then the value of $k$ is
\begin{qoptions}
\opt{A}{3}
\opt{B}{6}
\opt{C}{9}
\opt{D}{12}
\end{qoptions}
\qn{3} If $I_n=\displaystyle\int\cot^nx\,dx$, then $I_0+I_1+2(I_2+I_3+\ldots.I_8)+I_9+I_{10}$ equals;
\begin{qoptions}
\opt{A}{$\mu+\dfrac{\mu^2}2+\ldots+\dfrac{\mu^9}9+C$}
\opt{B}{$-\left\{\mu+\dfrac{\mu^2}2+\ldots+\dfrac{\mu^9}9\right\}+C$}
\opt{C}{$-\left\{\mu+\dfrac{\mu^2}{2!}+\ldots+\dfrac{\mu^9}{9!}\right\}+C$}
\opt{D}{$\dfrac\mu2+\dfrac{2\mu^2}3+\ldots+\dfrac{9\mu^9}{10}+C$}
\end{qoptions}
\{where $\mu=\cot x$\}
\qn{4} Integral $\displaystyle\int\dfrac{x^2+20}{(x\sin x+5\cos x)^2}dx.$ is
\begin{qoptions}
\opt{A}{$\dfrac{-x\sec x}{x\sin x+5\cos x}-\tan x+C$}
\opt{B}{$\dfrac{-x\sec x}{x\sin x+5\cos x}+\cot x+C$}
\opt{C}{$\dfrac{-x\sec x}{x\sin x+5\cos x}+\tan x+C$}
\opt{D}{$\dfrac{x\sec x}{x\sin x+5\cos x}-\tan x+C$}
\end{qoptions}
\qn{5} The parametric angle $\theta$, where $-\pi<\theta\le\pi$ of the point on the ellipse $\dfrac{x^2}{a^2}+\dfrac{y^2}{b^2}=1$ at which the tangent drawn cuts the intercept of minimum length on the co-ordinates axis, cannot be;
\begin{qoptions}
\opt{A}{$\tan^{-1}\sqrt{\dfrac ba}$}
\opt{B}{$-\tan^{-1}\sqrt{\dfrac ba}$}
\opt{C}{$\pi-\tan^{-1}\sqrt{\dfrac ba}$}
\opt{D}{$\pi+\tan^{-1}\left(\sqrt{\dfrac ab}\right)$}
\end{qoptions}
\qn{6} The asymptotes of the hyperbola $(x+y+1)^2-(x-y-3)^2=5$ cuts each other at $A$ and the coordinate axis at $B$ and $C$ then radius of circle passing through the points $A$, $B$, $C$ is
\begin{qoptions}
\opt{A}{3}
\opt{B}{$\dfrac{\sqrt5}2$}
\opt{C}{$\dfrac{\sqrt3}2$}
\opt{D}{$\sqrt2$}
\end{qoptions}
\sect{SECTION-I (ii)}
\qn{7} Consider a hyperbola $H$: $x^2-y^2=k$ and a parabola $P$ : $y=x^2$ then identify the correct statement(s)
\begin{qoptions}
\opt{A}{If point of intersections of $P$ and $H$ are concyclic then $k>2$}
\opt{B}{If $P$ and $H$ touch each other then $k=1/4$}
\opt{C}{If $k=-1/3$ and $m_1,m_2$ are the slopes of common tangents to $P$ and $H$ then $(3m_1^2+8)(3m_2^2+8)=112$}
\opt{D}{If $P$, $H$ do not touch but intersect at exactly two points then $k<0$}
\end{qoptions}
\qn{8} The coordinates of focus of a parabola which touches $x=1$, $y=1$, $x+2y=2$ and $y=2x+2$ are
\begin{qoptions}
\opt{A}{$\left(-\dfrac43,-\dfrac{13}3\right)$}
\opt{B}{$\left(-\dfrac{22}{31},-\dfrac73\right)$}
\opt{C}{$\left(-\dfrac1{20},\dfrac{17}{20}\right)$}
\opt{D}{None of these}
\end{qoptions}
\qn{9} If $\displaystyle\int\dfrac{dx}{x+x^{2011}}=f(x)+C_1$ and $\displaystyle\int\dfrac{x^{2009}}{1+x^{2010}}dx=g(x)+C_2$ (Where $C_1$ and $C_2$ are constants of integration). Let $h(x)=f(x)+g(x)$. If $h(1)=0$, then $h(e)$ is less than,
\begin{qoptions}
\opt{A}{0}
\opt{B}{1}
\opt{C}{$e$}
\opt{D}{2}
\end{qoptions}
\qn{10} A parabola $C$ whose focus is $S(0,0)$ and passing through $P(3,4)$. Equation to tangent at $P$ to parabola is $3x+4y-25=0$. A chord through $S$ parallel to tangent at $P$ intersects the parabola at $A$ and $B$, then
\begin{qoptions}
\opt{A}{Length of $AB$ is 20 units}
\opt{B}{Area of triangle $APS$ is 40 square units}
\opt{C}{Latus rectum of parabola is 20 units}
\opt{D}{Area of triangle $PAB$ is 80 square units}
\end{qoptions}
\qn{11} If {\small\[\int(x^{2010}+x^{804}+x^{402})(2x^{1608}+5x^{402}+10)^{\frac1{402}}dx\]\[=\dfrac1{10a}(2x^{2010}+5x^{804}+10x^{402})^{\frac a{402}}+C\]} then $a$ is less than
\begin{qoptions}
\opt{A}{401}
\opt{B}{403}
\opt{C}{411}
\opt{D}{413}
\end{qoptions}
\qn{12} The circle centered at origin is inscribed in the parabola $y=x^2-100$. Which of the following is/are true?
\begin{qoptions}
\opt{A}{Radius of circle is $\dfrac{\sqrt{399}}2$}
\opt{B}{Radius of circle is $\sqrt{\dfrac{399}2}$}
\opt{C}{Common tangent to circle and parabola passes from $\left(0,\dfrac{-399}2\right)$}
\opt{D}{Length of common chord for circle and parabola is $\dfrac{\sqrt{398}}2$}
\end{qoptions}
\sect{SECTION-II}
\qn{13} If {\small\[\int(\cot2x\cot3x-\tan2x\tan7x)\tan5x\,dx\]\[=a\,\ell n(\tan2x)+b\,\ell n(\sin3x)\]\[\qquad+c\,\ell n(\sec5x)+d\,\ell n(\cos7x)+C\]} and $a,b,c,d\in Q$ and $C$ is constant of integration. If $a+b+c+d$ can be expressed as $\dfrac mn$ in lowest form, then $m+n$ is $(m,n\in N)$
\qn{14} Variable pairs of chords at right angles are drawn through a point $P$ (with eccentric angle $\pi/4$) on the ellipse $\dfrac{x^2}4+y^2=1$, to meet the ellipse at two point say $A$ and $B$. If the line joining $A$ and $B$ passes through a fixed point $Q(a,b)$ such that $a^2+b^2$ has the value equal to $\dfrac mn$, where $m$, $n$ are relatively prime positive integers, find $(n-m)$.
\qn{15} An equilateral triangle $ABC$ is inscribed in the parabola $y=x^2$ and one of the side of the equilateral triangle has the gradient 2. If the sum of $x$-coordinates of the vertices of the triangle is a rational in the form $p/q$ where $p$ and $q$ are coprime, then find the value of $\dfrac{(p+q)}2$
\qn{16} $C$ is the centre of the hyperbola $\dfrac{x^2}4-\dfrac{y^2}1=1$ and $A$ is any point on it. The tangents at $A$ to the hyperbola meets the line $x-2y=0$ and $x+2y=0$ at $Q$ and $R$ respectively then $CQ.CR=$
\qn{17} Let $\displaystyle\int\sec^{-1}[-\sin^2x]\,dx=f(x)+C$ where $[.]$ denotes greater integer function, then the value of {\small$\left(f\left(\dfrac8{\pi x}\right)\right)''$} at $x=2$, is
\qn{18} At a point $A(1,1)$ on ellipse, equation of tangent is $y=x$. If one of the foci of the ellipse is $(0,-2)$ and the co-ordinates of centre of ellipse are $(\alpha,\beta)$ the value of $\alpha+\beta$ is (given length of major axis of ellipse is $4\sqrt{10}$ units)
\end{multicols}
\end{document}
